\documentclass[11pt]{article}
\usepackage[margin=1in]{geometry}
\usepackage{amsmath,amssymb}
\usepackage{booktabs}
\usepackage{longtable}
\usepackage{graphicx}
\usepackage{microtype}
\usepackage[hidelinks]{hyperref}
\usepackage{siunitx}
\title{Connectivity-defined coordinates and dynamical motifs in multitask
recurrent networks: methods and results}
\author{Working report}
\date{2026-07-27}

% Factual report. States methods and numerical results with few conclusions. Every
% comparison is reported against its size-matched null; negative and null results are
% reported alongside positive ones. Primary results are reported across a k-sweep rather than
% at a single operating scale. The focus is the ten corrected 256-unit networks; other sizes
% appear only through their released seeds 0-2. The uncorrected first-batch retrained networks
% and their superseded outputs are documented separately in legacy_uncorrected_seeds_report.

\begin{document}
\maketitle

\begin{abstract}
This report documents the methods and numerical results of an analysis of trained multitask
recurrent neural networks (RNNs). We identify approximate equitable partitions of the recurrent
weight matrix, quantify how far the resulting group-averaged (``quotient'') coordinates capture
the recurrent dynamics, locate the networks' fixed-point motifs across task epochs, and relate
the motif directions to the connectivity coordinates. The principal model set is ten $256$-unit
networks (three released, seven trained for this study); other network sizes appear only through
their released seeds. Results are reported across a sweep of the group number $k$ rather than at
a single operating scale. Every comparison is reported against its size-matched null and, where
computed, against an activity-defined partition. Positive, negative, and null outcomes are
reported with equal prominence. The unit of replication is the trained network (seed).
\end{abstract}

\tableofcontents

\section{Model set and training}

\subsection{Architecture and task set}
We analysed leaky discrete-time RNNs with softplus activation $\sigma(x)=\log(1+e^{x})$ and
diagonal recurrent initialisation, matching the primary model class of Driscoll, Shenoy \&
Sussillo (2024). The recurrent state $h(t)\in\mathbb{R}^{N}$ evolves as
\[
h(t{+}1)=(1-\alpha)\,h(t)+\alpha\,\sigma\!\left(W_{\rm rec}h(t)+W_{\rm in}u(t)+b_{\rm rec}\right),
\qquad z(t)=W_{\rm out}h(t)+b_{\rm out},
\]
with leak $\alpha=0.2$ (recovered exactly from stored Jacobians; \S\ref{sec:quotient}). Inputs
$u(t)$ comprise one fixation channel, four stimulus channels (two circular modalities as
sine/cosine pairs), and $15$ one-hot rule channels (input dimension $20$); outputs comprise one
fixation and two response-direction channels. The task set is the $15$ cognitive tasks of the
multitask paradigm. All hidden-activity analyses used noiseless test trials.

\subsection{Networks analysed}
\label{sec:nets}
The principal analyses use ten independently trained $256$-unit networks (seeds $0$--$9$): seeds
$0$--$2$ are the networks released by Driscoll et al.; seeds $3$--$9$ were trained for this study
on the corrected objective (below) and have L1 activity/weight coefficients $0$, L2 coefficients
$10^{-6}$, and $212$--$220$ active units (task-epoch activity variance $>10^{-3}$). Mean accuracy
across tasks is $\approx0.97$--$0.99$ (released) and $\approx0.97$ (retrained; per-task range:
non-decision tasks $0.998$--$1.000$, decision-making tasks $0.94$--$0.96$). Networks at $128$,
$512$ and $1024$ units are represented only by their released seeds $0$--$2$ and are used solely
for size comparisons. An earlier batch of retrained seeds was produced under an unintended
regularisation coefficient and has been superseded; that batch and its outputs are documented in
the companion report \texttt{legacy\_uncorrected\_seeds\_report}.

\section{Structural partition identification}
\label{sec:partition}
Analyses were restricted to active units. The recurrent block was split into
$W^{+}=\max(W,0)$, $W^{-}=\max(-W,0)$ with $W=W_{\rm rec}^{T}$ (each unit characterised by the
signed pattern of recurrent inputs it receives). A partition into $k$ groups is an indicator
matrix $H\in\{0,1\}^{n\times k}$; equitability is scored by
\[
\Gamma_{\rm signed}(H)=\lVert W^{+}H-HD^{-1}H^{T}W^{+}H\rVert_{F}
+\lVert W^{-}H-HD^{-1}H^{T}W^{-}H\rVert_{F},\qquad D=H^{T}H,
\]
minimised by an approximate signed Weisfeiler--Leman refinement ($k$-means reclustering of the
concatenated profile $[W^{+}H\,|\,W^{-}H]$; $20$ iterations, $20$ initialisations). The full
$k$-sweep ($k=3$--$60$) was cached, together with an independently computed activity partition
(Ward clustering of normalised task-epoch variance profiles) at each $k$.

\paragraph{Reporting scale.} Results below are reported across a sweep $k\in\{8,12,16,20,25,30,
40,50\}$. A single ``operating'' scale of $20\%$ of the active-unit count (giving $k=33$--$51$
for the $256$-unit networks) is used only where one scale is required; it is not the basis of any
conclusion, and the $k$-sweep is the primary presentation.

\paragraph{Group-size distribution and singletons.} The refinement does not constrain group
sizes; groups are unequal and some are singletons. Median over the ten $256$-unit networks of the
connectivity partition (of $\approx210$ active units): at $k=8$, group sizes $1/26/63$
(min/median/max) with $1$ singleton; at $k=20$, $1/9/26$ with $2$ singletons; at $k=30$,
$1/6/25$ with $5$ singletons; at $k=50$, $1/3/18$ with $14$ singletons. Singleton groups are
retained (they carry a one-unit row of the quotient); no size floor is imposed. Because every
comparison partition (random, activity, spectral, modularity) is matched to the same group-size
multiset (\S\ref{sec:statconv}), the effect of the size distribution is controlled rather than
assumed away.

\section{Lumpability of the linear recurrent map (Result 1)}
\label{sec:lump}
Let $P=H(H^{T}H)^{-1}H^{T}$. For each leading eigenvector $v$ of a matrix $M$ we measured the
leakage $\lVert v-Pv\rVert/\lVert v\rVert$, step $R^{2}=1-\lVert v-Pv\rVert^{2}/\lVert v-\bar
v\rVert^{2}$, and the subspace invariance $\lVert(I-P)MP\rVert/\lVert MP\rVert$, each
standardised against size-matched random partitions ($200$ permutations) and reported as a
signed \emph{score} (positive-for-more-structure). $M=W_{\rm rec}^{T}$ and the raw state-update
matrix were both evaluated.

\subsection{Results (k-sweep, ten $256$-unit networks)}
\begin{center}\small
\begin{tabular}{rrrrrr}
\toprule
$k$ & leak score ($W^{T}$) & \#$>0$ & leak score ($W$) & invariance score & \#$>0$ (inv) \\
\midrule
8  & +27.8 & 10/10 & +21.9 & +1.1  & 7/10 \\
12 & +29.8 & 10/10 & +22.6 & +2.8  & 7/10 \\
16 & +28.4 & 10/10 & +25.3 & +5.1  & 8/10 \\
20 & +27.8 & 10/10 & +21.2 & +7.0  & 9/10 \\
25 & +22.3 & 10/10 & +20.0 & +9.4  & 9/10 \\
30 & +21.2 & 10/10 & +18.1 & +11.7 & 10/10 \\
40 & +20.3 & 10/10 & +14.7 & +16.8 & 10/10 \\
50 & +19.4 & 10/10 & +12.8 & +17.7 & 10/10 \\
\bottomrule
\end{tabular}
\end{center}
The eigenvector-leakage score was positive in $10/10$ networks at every $k$; the
subspace-invariance score rose with $k$ and was positive in $7$--$10/10$ (weaker at small $k$).
For the canonical network (seed $1$, $k=51$) a leading eigenvector had step $R^{2}=0.89$ on the
recovered groups versus $R^{2}=0.20$ under a size-matched random regrouping. For the released
seeds $0$--$2$ at other sizes (operating scale), the leakage score was $+12.9$ ($128$), $+22.6$
($512$), $+26.8$ ($1024$).

\subsection{Degree/training surrogate control}
\label{sec:surrogate}
The identical pipeline (WL optimisation at the operating $k$, then equitability $z$ and step
$R^{2}$) was applied to three structure-destroying surrogates: (i) a row shuffle preserving each
unit's incoming-weight strength; (ii) a full shuffle; (iii) an i.i.d.\ Gaussian matched in
Frobenius norm. For the ten corrected $256$-unit networks:
\begin{center}\small
\begin{tabular}{lcc}
\toprule
matrix & equitability score & step $R^{2}$ \\
\midrule
trained            & $+7.35$  & $0.784$ \\
row shuffle        & $+12.26$ & $0.693$ \\
full shuffle       & $+3.14$  & $0.415$ \\
matched Gaussian   & $+2.12$  & $0.275$ \\
\bottomrule
\end{tabular}
\end{center}
The trained matrix exceeded the row shuffle on step $R^{2}$ in $10/10$ networks and on
equitability score in $0/10$ (the row shuffle is more equitable than the trained matrix in every
network). Equitability magnitude is thus not higher for the trained matrix than for a
strength-preserving shuffle; step $R^{2}$ is.

The reason the row shuffle scores \emph{higher} on equitability is a property of the null, not of
trained structure. The equitability score is $z=(\overline{\Gamma_{\rm null}}-\Gamma_{\rm obs})/
\mathrm{sd}(\Gamma_{\rm null})$, where $\Gamma_{\rm null}$ is recomputed from each surrogate's own
size-matched random partitions on that surrogate's own matrix (each surrogate is scored against
its own null, verified in code). On a representative network (seed $0$, $k=20$): the trained
matrix gave $\Gamma_{\rm obs}=157.5$ against its own null mean $151.8$ ($\mathrm{sd}\ 2.5$),
i.e.\ $z=-2.25$ --- the WL partition of the trained matrix is \emph{not} more equitable than a
random same-size partition of it. The row-shuffled matrix gave $\Gamma_{\rm obs}=128.0$ against
its own null mean $161.9$ ($\mathrm{sd}\ 1.7$), i.e.\ $z=+19.5$. Within-row shuffling makes rows
exchangeable and separable essentially by incoming strength (degree); WL can then group rows to
drive $\Gamma$ far below a degree-blind random null, whereas the trained matrix's structure is
not of that degree-separable kind. The equitability $z$ thus rewards degree separability, which
the shuffle amplifies; this is why the load-bearing structural claim is lumpability (training-made,
trained $>$ shuffle $10/10$; \S\ref{sec:lump}) rather than equitability magnitude.

\section{Partition non-identifiability and subspace stability}
\label{sec:nonident}
Re-running the refinement from $20$ initialisations (256-unit networks): the cost
$\Gamma_{\rm signed}$ was reproducible (coefficient of variation $\approx3\%$) while the specific
assignments were not (mean pairwise adjusted Rand index $\approx0.40$). The group-constant
subspaces of restart pairs had mean squared principal-angle cosine $\approx0.63$ versus
$\approx0.24$ for size-matched random partitions (mean principal angle $\approx34^{\circ}$).
Across-restart coefficients of variation were $\approx2.5\%$ (step-$R^{2}$) and $\approx10\%$
(leakage). Single-partition analyses used a co-association consensus.

\section{Comparison to spectral and modularity partitions}
\label{sec:partfamily}
A control for whether the backbone results require the signed-WL partition specifically, or would
follow from any sensible graph partition at matched $k$. At each $k$ we built, on the same active
recurrent block, a normalised spectral-clustering partition (affinity $|W|+|W|^{T}$) and a greedy
modularity partition (same weighted graph, forced to $k$ communities), and evaluated the three
label-vector metrics behind Results 1/3/4: eigenvector leakage (Result 1, lower is more lumpable),
median motif-direction containment (Result 3, higher), and top-3 quotient-mode energy (Result 4,
higher). Medians over the ten $256$-unit networks; signed-rank compares WL to each alternative.
\begin{center}\small
\begin{tabular}{lrccccc}
\toprule
metric & $k$ & WL & spectral & modularity & random & WL vs spec / mod ($p$) \\
\midrule
leakage $\downarrow$     & 12 & 0.656 & 0.806 & 0.856 & 0.868 & 0.002 / 0.002 \\
                         & 20 & 0.644 & 0.795 & 0.855 & 0.863 & 0.002 / 0.002 \\
                         & 30 & 0.635 & 0.751 & 0.842 & 0.830 & 0.002 / 0.002 \\
containment $\uparrow$   & 12 & 0.470 & 0.222 & 0.124 & 0.222 & 0.002 / 0.002 \\
                         & 20 & 0.598 & 0.285 & 0.166 & 0.294 & 0.002 / 0.002 \\
                         & 30 & 0.677 & 0.367 & 0.245 & 0.365 & 0.002 / 0.002 \\
top-3 $\uparrow$         & 12 & 0.802 & 0.511 & 0.623 & 0.547 & 0.002 / 0.002 \\
                         & 20 & 0.641 & 0.360 & 0.503 & 0.426 & 0.002 / 0.004 \\
                         & 30 & 0.560 & 0.318 & 0.448 & 0.348 & 0.002 / 0.002 \\
\bottomrule
\end{tabular}
\end{center}
The signed-WL partition exceeded both spectral and modularity partitions on all three metrics at
every $k$ ($10/10$ networks, $p\leq0.004$). The spectral and modularity partitions were at or
below the size-matched random floor on containment (modularity below random, since dense unsigned
communities are not aligned with the motif subspace). The lumpable/containing/concentrating
coordinates therefore come specifically from the signed-equitability partition, not from graph
community structure in general. (The leakage advantage is partly definitional --- WL minimises an
equitability cost tied to eigenvector leakage --- but containment and top-3 energy are motif
quantities that WL does not optimise, and it wins there too.)

\section{Fixed-point motifs}
\label{sec:motifs}
For each task epoch, holding the epoch input constant, fixed and slow points of the autonomous
dynamics were located by gradient descent on the velocity
$q(h)=\tfrac12\lVert -h+\sigma(W_{\rm rec}h+W_{\rm in}u+b)\rVert^{2}$ (thresholds $q<10^{-8}$ for
isolated fixed points, $q<10^{-6}$ for the slow cloud). Analysed epochs were fixation, stimulus
(\texttt{stim1/2}) and delay (\texttt{delay1/2}); the pure response epoch (\texttt{go1}) was
analysed separately (\S\ref{sec:rotation}). Each epoch was classified from local Jacobian
spectra as a point attractor (all $|\lambda|<1$), a saddle ($\geq1$ mode $|\lambda|>1$), or a
line/ring attractor (marginal mode, locally one-dimensional cloud); a rotational motif would show
a Jacobian eigenvalue with imaginary part $>0.05$.

\subsection{Taxonomy (corrected models, all ten seeds)}
Across $490$ classified task epochs the inventory was $263$ line attractors, $101$ point
attractors, $76$ saddles, $50$ mixed, and $0$ rotations; all fixed-point Jacobian spectra were
near-real. Epoch-type gate counts (matched $/$ total): pooled over seeds $0$--$2$, memory delays
$\to$ line attractor $40/48$, decision stimulus $\to\geq1$ saddle $30/42$, sensory stimulus
$\to$ point attractor $9/12$; pooled over all ten seeds, $130/160$, $109/140$, and $13/40$. Each
motif was summarised by a unit direction $p$: the marginal (memory) mode of a line attractor or
the leading unstable (decision) axis of a saddle.

\section{Motif directions in the connectivity subspace (Result 3, containment)}
\label{sec:containment}
We measured the fraction of each motif direction inside the group-constant subspace,
$\lVert Pp\rVert/\lVert p\rVert$, and at matched $k$ the excess containment
$\Delta=f_{\rm obs}-\overline{f_{\rm rand}}$ for the connectivity and activity partitions (each
minus its own size baseline). Across the ten $256$-unit networks:
\begin{center}\small
\begin{tabular}{rrrrr}
\toprule
$k$ & $\Delta_{\rm role}$ & $\Delta_{\rm activity}$ & $z$ (role$-$act) & role $>$ act \\
\midrule
8  & 0.121 & 0.016 & +4.01 & 10/10 \\
12 & 0.183 & 0.010 & +4.97 & 10/10 \\
16 & 0.239 & 0.014 & +6.35 & 10/10 \\
20 & 0.280 & 0.012 & +7.31 & 10/10 \\
25 & 0.295 & 0.021 & +6.24 & 10/10 \\
30 & 0.282 & 0.020 & +6.04 & 10/10 \\
40 & 0.268 & 0.018 & +6.03 & 10/10 \\
50 & 0.270 & 0.037 & +5.13 & 10/10 \\
\bottomrule
\end{tabular}
\end{center}
The connectivity excess containment exceeded the activity partition's in $10/10$ networks at
every $k$; the activity excess was near zero throughout.

\subsection{$W_{\rm rec}$-eigenbasis null (strong null; negative for containment)}
\label{sec:wrec-null}
The random and activity partitions are weak nulls. A stronger, connectivity-derived null is the
eigenbasis of $W_{\rm rec}$ itself: the span of the top-$k$ eigenvectors of $W_{\rm rec}$ (by
$\lvert\lambda\rvert$, real parts, orthonormalised) matched in dimension to the $k$-group
constant subspace. We compared the raw containment $\lVert Pp\rVert/\lVert p\rVert$ of the motif
directions onto each subspace (median over the ten networks; two-sided signed-rank):
\begin{center}\small
\begin{tabular}{rcccc}
\toprule
$k$ & containment (conn) & containment ($W_{\rm rec}$-eig) & conn $>$ eig (motif frac.) & Wilcoxon $p$ \\
\midrule
8  & 0.320 & 0.678 & 0.00 & 0.002 \\
12 & 0.470 & 0.706 & 0.00 & 0.002 \\
16 & 0.540 & 0.716 & 0.02 & 0.002 \\
20 & 0.598 & 0.729 & 0.03 & 0.002 \\
25 & 0.652 & 0.743 & 0.04 & 0.004 \\
30 & 0.677 & 0.757 & 0.08 & 0.002 \\
40 & 0.716 & 0.776 & 0.19 & 0.014 \\
50 & 0.769 & 0.797 & 0.23 & 0.160 \\
\bottomrule
\end{tabular}
\end{center}
The top-$k$ $W_{\rm rec}$ eigenbasis contained the motif directions more than the connectivity
group-constant subspace at every $k$ (signed-rank in the eigenbasis's favour up to $k=40$;
per-motif the connectivity subspace won in $0$--$23\%$ of motifs). The connectivity partition is
therefore \emph{not} a privileged $k$-dimensional subspace for the motif directions relative to
the leading eigenvectors of the same matrix; its containment advantage in
\S\ref{sec:containment} is specifically over the activity partition and size-matched random
partitions, not over the $W_{\rm rec}$ spectrum. This is expected --- the leading eigenvectors
are by construction the best $k$-dimensional linear fit to the slow directions the motifs align
with --- but it bounds the containment claim and is stated because it is the first null a
reader is likely to name.

A second natural $k$-dimensional subspace, defined by the task inputs rather than by recurrence,
is the span of the top-$k$ right-singular directions of $W_{\rm in}$ over the active units. It
too contained the motif directions more than the connectivity groups at every $k$ (median
containment $0.75$ vs.\ $0.60$ at $k=20$; connectivity above the input span in $1$--$13\%$ of
motifs; signed-rank $p=0.002$ at all $k$). The connectivity partition's containment advantage is
therefore specifically over the \emph{activity} and size-matched \emph{random} partitions
(\S\ref{sec:containment}); it is not an advantage over either the leading recurrent eigenvectors
or the input-weight span.

\section{Quotient reduced model (Result 2)}
\label{sec:quotient}
At a fixed point $J=(1-\alpha)I+\alpha\,\mathrm{diag}(g)\,W_{\rm rec}$ with $g=\sigma'(x^{*})$;
$\alpha=0.200$ and the per-unit gains $g$ are recovered exactly ($\sim10^{-8}$) from a stored
Jacobian and $W_{\rm rec}$, with no fitting. We formed $Q_{\rm struct}=D^{-1}H^{T}W_{\rm rec}H$,
$Q_{\rm pred}=(1-\alpha)I+\alpha\,\mathrm{diag}(\bar g)\,Q_{\rm struct}$, and $Q_{\rm task}=
D^{-1}H^{T}\bar J H$. The top-$m$ eigenvectors of $Q$, lifted to units and orthonormalised, form
$B$; for a motif direction $v$, $A_{\rm raw}=\lVert B^{T}v\rVert/\lVert v\rVert$ and
$A_{\rm cond}=\lVert B^{T}v\rVert/\lVert P_{H}v\rVert$. Nulls are size-matched random partitions
with their own recomputed quotient, plus the activity partition at matched $k$.

\subsection{Prediction across the k-sweep ($Q_{\rm pred}$, $r=3$, ten networks)}
\begin{center}\small
\begin{tabular}{rrrrr}
\toprule
$k$ & $A_{\rm cond}$ (conn) & random & conn $>$ rand & $z$ \\
\midrule
8  & 0.932 & 0.736 & 10/10 & +2.62 \\
12 & 0.890 & 0.660 & 10/10 & +3.56 \\
16 & 0.874 & 0.625 & 10/10 & +3.88 \\
20 & 0.816 & 0.612 & 10/10 & +3.45 \\
25 & 0.809 & 0.598 & 10/10 & +3.32 \\
30 & 0.797 & 0.595 & 10/10 & +3.42 \\
40 & 0.783 & 0.594 & 10/10 & +3.37 \\
50 & 0.787 & 0.611 & 10/10 & +3.07 \\
\bottomrule
\end{tabular}
\end{center}
Connectivity exceeded its random-quotient null in $10/10$ networks at every $k$. Eigenvalue
agreement at $k=20$: $|\lambda_{Q}-\lambda_{J}|=0.046$ (connectivity) vs.\ $0.191$ (random) and
$0.141$ (activity).

\subsection{Compression ($A_{\rm raw}$, $k=20$, $r=3$, median over ten networks)}
Quotient $Q_{\rm pred}$ ($k^{2}+k=420$ free parameters) $=0.479$; full mean Jacobian
($\approx4\times10^{4}$ parameters) $=0.801$; full $W_{\rm rec}$ top-$3$ eigenbasis $=0.466$;
$W_{\rm rec}$ SVD basis $=0.113$; random directions in the group subspace $=0.205$; containment
ceiling $\lVert P_{H}p\rVert=0.578$.

\subsection{Gain ablations ($A_{\rm cond}$, $k=20$, $r=3$, median over ten networks)}
\begin{center}\small
\begin{tabular}{lc}
\toprule
predictor & $A_{\rm cond}$ \\
\midrule
$Q_{\rm pred}$ (group-mean gains)          & $0.816$ \\
scalar gain (one global)                    & $0.749$ \\
permuted group gains                        & $0.680$ \\
no gains ($Q_{\rm struct}$)                 & $0.670$ \\
cross-condition gains                       & $0.775$ \\
gains only (no quotient)                    & $0.460$ \\
size-matched random (null)                  & $0.621$ \\
\bottomrule
\end{tabular}
\end{center}
The gains-only predictor falls below the random-partition null; the no-gain quotient exceeds it
by $0.049$.

\subsection{Eigenvalue-matched non-motif control (negative)}
\label{sec:eigmatched}
$Q_{\rm pred}$'s prediction of the motif direction $v$ was compared with its prediction of a
non-motif eigenvector $q$ of the \emph{same} mean Jacobian, chosen with $|\lambda_{q}|$ nearest
$|\lambda_{v}|$ and alignment to $v<0.30$; partition, gains, $r$ and conditioning identical.
Across the ten networks the fraction of conditions in which the motif direction was predicted
better than the eigenvalue-matched non-motif direction was: $k=8$, $54\%$; $k=12$, $52\%$;
$k=16$, $50\%$; $k=20$, $46\%$; $k=25$, $50\%$; $k=30$, $49\%$; $k=40$, $48\%$; $k=50$, $43\%$
(of $440$ conditions). At $k=20$ the per-network gap ($A_{\rm cond}$ motif $-$ matched) was
$+0.026$ (median), positive in $7/10$ networks; by motif type the per-condition win rate was
line $138/263$, saddle $32/76$, point $34/101$.

\subsection{Held-out gains (prospective, seeds $0$--$2$)}
Estimating gains on one contiguous ring arc and predicting tangent directions at held-out angles
gave $A_{\rm cond}=0.736$ vs.\ random $0.597$ ($k=20$); interleaved held-out $\approx$ in-sample
($0.866$ vs.\ $0.867$); cross-task gains $0.812$, cross-type gains $0.765$.

\section{Task accessibility of motif vs.\ matched non-motif modes (delay epoch; negative)}
\label{sec:access}
For each motif direction $v$ and its eigenvalue-matched non-motif direction $q$
(\S\ref{sec:eigmatched}), evaluated at the fixed-point epoch, seven task-interface measures were
compared (seeds $0$--$2$, $128$ conditions):
\begin{center}\small
\begin{tabular}{lccc}
\toprule
measure & motif & matched & motif $>$ matched \\
\midrule
input-drive projection $|\cos(m,Bu)|$        & 0.076 & 0.123 & 45/128 \\
state-displacement projection                 & 0.160 & 0.208 & 62/128 \\
trajectory variance along mode                & 0.243 & 0.128 & 86/128 \\
readout potency $\lVert R^{T}m\rVert^{2}$     & 0.001 & 0.001 & 53/128 \\
persistence (local $|\lambda|$)               & 0.975 & 0.970 & 75/128 \\
modal controllability $\lVert B^{T}w\rVert$   & 0.030 & 0.035 & 45/128 \\
modal observability $\lVert C\,m\rVert$       & 0.101 & 0.121 & 55/128 \\
\bottomrule
\end{tabular}
\end{center}
$B=\alpha\,\mathrm{diag}(g)\,W_{\rm in}$ is the input Jacobian, $C$ the response readout, $R$ its
orthonormal basis. Readout potency was $\approx0.001$ for both modes (chance $\approx0.010$).

\section{Input-interpolation branch (re-found manifold; negative)}
\label{sec:branch}
This is a negative result: the reduced quotient system does not reproduce the full network's
re-found slow manifold as the input is interpolated. For a memory-ring task the input
$u(\beta)=(1-\beta)u_{\rm delay}+\beta u_{\rm response}$ ramps the fixation channel from $1$ to
$0$. At each $\beta$ the slow set was re-found by relaxing a broad seed pool under the leaky
dynamics; the full network and the reduced quotient system were compared. For \texttt{delaygo},
canonical network ($k=20$): the full system had a one-dimensional slow set up to
$\beta\approx0.9$ while the reduced system had a single fixed point (effective dimension $0$)
throughout; marginal eigenvalue full $0.95$--$0.999$, reduced $0.96$--$0.98$; centroid cosine
$0.74$--$0.98$; slow-subspace principal angle $43^{\circ}$--$63^{\circ}$; readout potency rose
$0.000\to0.092$ (full) and $0.006\to0.026$ (reduced).

\paragraph{Sweep over $k$, task and seed.} To show the failure is not specific to one setting the
comparison was repeated over two tasks (\texttt{delaygo}, \texttt{delayanti}), two seeds ($0$
released, $3$ corrected) and $k\in\{12,20,30\}$, averaging each metric over $\beta$.
\begin{center}\small
\begin{tabular}{llcccc}
\toprule
seed & task & eff.\ dim (full) & eff.\ dim (reduced) & slow-plane angle & potency full / reduced \\
\midrule
0 & \texttt{delaygo}   & 0.73 & 0.00 & $52$--$58^{\circ}$ & 0.024 / 0.015 \\
0 & \texttt{delayanti} & 1.20 & 0.00 & $50$--$53^{\circ}$ & 0.033 / 0.011 \\
3 & \texttt{delaygo}   & 1.64 & 0.00 & $73$--$78^{\circ}$ & 0.008 / 0.013 \\
3 & \texttt{delayanti} & 2.32 & 0.00 & $77$--$79^{\circ}$ & 0.009 / 0.013 \\
\bottomrule
\end{tabular}
\end{center}
Values are averaged over $\beta$; the angle range spans $k\in\{12,20,30\}$. In all $12$ cells the
reduced effective dimension was $0$ against a full effective dimension $>0.7$.
The reduced system's slow set collapsed (effective dimension near $0$) while the full system's
persisted, and the reduced slow plane stayed far from the full one, in every cell. The quotient
captures the coarse existence and rough location of the memory manifold but not its dimensionality
or fine geometry along the input interpolation, at any of the tested scales.

\section{Rotation into the readout at the response epoch (Result 3, decisive control)}
\label{sec:rotation}
At the response epoch we compared the readout potency of the response-epoch ring plane
(top-$2$ PCs of the response fixed-point cloud) with a matched non-ring plane (two eigenvectors
of the response mean Jacobian with $|\lambda|$ nearest the ring modes', alignment $<0.40$) and a
random plane. Across the ten networks ($109$ delayed-response tasks), median readout potency:
ring at the delay epoch $0.0004$; ring at the response epoch $0.0302$; matched non-ring
$0.0015$; random plane $0.0086$. The ring plane exceeded the matched non-ring plane in $104/109$
tasks (released seeds $33/33$; retrained seeds $71/76$) and exceeded the random plane in
$100/109$; per network the median ring potency exceeded the median matched potency in $10/10$.

\section{Concentration of motifs on quotient eigenmodes (Result 4)}
\label{sec:concentration}
This is a \emph{within-network} statement: for each network separately, the collection of that
network's task-epoch motifs is expressed in the coordinates of that network's own quotient, and we
ask whether the motifs load a few of those modes rather than spreading over all $k$. No claim is
made that a particular mode index recurs across networks --- mode $j$ of one network is not mode
$j$ of another --- so all statistics are computed per network and only the resulting scalar is
aggregated across seeds. The primary statistic is the fraction of aggregate motif energy on the
top three modes; the secondary statistic is the participation ratio
$\mathrm{PR}=(\sum_{j}E_{j})^{2}/\sum_{j}E_{j}^{2}$ (effective number of modes). Because the
quotient is directed and non-normal its eigenvectors are not orthogonal, so the modes are
orthonormalised (the lifted eigenvectors are Gram--Schmidt/QR ordered by $\lvert\lambda\rvert$)
before energy is measured; \S\ref{sec:typeprofiles} confirms that the low PR is not an artifact of
this choice (a complex-aware orthonormal basis and the real Schur vectors give the same PR). The
null is size-matched random partitions with their own quotient modes; significance is a two-sided
Wilcoxon signed-rank test across the ten networks.
\begin{center}\small
\begin{tabular}{rccccc}
\toprule
$k$ & top-3 (conn) & top-3 (random) & conn $>$ rand & Wilcoxon $p$ & PR conn / rand \\
\midrule
12 & 0.802 & 0.546 & 10/10 & 0.0020 & 2.72 / 7.03 \\
20 & 0.641 & 0.430 & 10/10 & 0.0020 & 4.87 / 10.55 \\
30 & 0.560 & 0.352 & 10/10 & 0.0020 & 6.59 / 14.70 \\
\bottomrule
\end{tabular}
\end{center}
The top-3 connectivity energy exceeded the random-quotient value in $10/10$ networks at every
$k$ ($p=0.0020$, the signed-rank floor for $n=10$). Correlation of the top-3 connectivity energy
with per-network task accuracy ($k=20$, $n=10$): Pearson $r=+0.47$ ($p=0.17$), Spearman
$\rho=+0.12$ ($p=0.75$); the participation ratio versus accuracy: Pearson $r=-0.46$ ($p=0.18$).
Neither correlation is significant at $n=10$.

\subsection{$W_{\rm rec}$-eigenbasis null (strong null; mixed)}
Against the same $W_{\rm rec}$-eigenbasis null used for containment (\S\ref{sec:wrec-null}), the
top-3 energy comparison is $k$-dependent rather than uniformly in favour of the connectivity
quotient. Median top-3 energy on the connectivity quotient's top-3 eigenmodes versus the top-3
$W_{\rm rec}$ eigenvectors, over the ten networks:
\begin{center}\small
\begin{tabular}{rccc}
\toprule
$k$ & top-3 (conn quotient) & top-3 ($W_{\rm rec}$-eig) & Wilcoxon $p$ \\
\midrule
8  & 0.888 & 0.779 & 0.105 \\
12 & 0.802 & 0.727 & 0.275 \\
16 & 0.757 & 0.707 & 0.846 \\
20 & 0.641 & 0.695 & 0.027 \\
30 & 0.560 & 0.646 & 0.002 \\
50 & 0.518 & 0.585 & 0.002 \\
\bottomrule
\end{tabular}
\end{center}
At small $k$ ($k\leq12$) the connectivity quotient concentrated motif energy on slightly fewer
modes than the $W_{\rm rec}$ eigenbasis, but not significantly; at $k\geq20$ the $W_{\rm rec}$
eigenbasis concentrated more (significantly from $k=30$). The concentration in
\S\ref{sec:concentration} is thus specific relative to random quotients but is not sharper than
the leading $W_{\rm rec}$ eigenvectors except at small $k$.

\section{Mode-energy profiles by motif type}
\label{sec:typeprofiles}
The concentration above pools all motif types. Here we ask, within each network, whether the three
motif types occupy the \emph{same} quotient modes or different ones, and we confirm the
concentration is robust to the coordinate system. For each network we formed three aggregate
mode-energy vectors --- line attractors, point attractors, saddles --- each an equal-weight mean of
the per-motif energy profiles (every motif's profile normalised to sum one \emph{before} pooling,
so that the numerous line-attractor samples do not dominate; a per-task-family average gives the
same conclusions). Energy is measured in an orthonormal basis of the quotient (below). All
comparisons are within a network; only the resulting scalars are aggregated across the ten seeds.

\paragraph{Coordinate system (non-normality check).} The quotient is directed and non-normal, so
its eigenvectors are not orthogonal and raw-eigenvector energy/PR could be distorted. We therefore
used two orthonormal bases whose leading columns span the dominant ($\lvert\lambda\rvert$-ordered)
invariant subspace: (i) an eigen-ordered basis retaining both the real and imaginary parts of each
complex conjugate pair (its full $2$-D real invariant plane), Gram--Schmidt orthonormalised; and
(ii) the real Schur vectors. The median per-motif participation ratio ($k=20$) was $3.67$ in the
eigen-ordered orthonormal basis and $3.54$ in the Schur basis, versus $3.33$ for the earlier
real-parts-only construction --- i.e.\ the proper orthonormal bases give, if anything, a
\emph{higher} PR. The low participation ratio is therefore not an artifact of non-orthogonal
eigenvectors. Within-type concentration held for all three types on their own network's modes
(equal-weighted top-3 energy: saddles $0.72$, lines $0.68$, points $0.51$; points least
concentrated).

\paragraph{Do the types share modes?} We compared the three type profiles pairwise by cosine and
correlation, by top-3 mode-index overlap, and by the principal angles between the top-3 mode spans,
each against size-matched random quotients (median over ten seeds; two-sided signed-rank of
observed vs.\ null).
\begin{center}\small
\begin{tabular}{lcccccc}
\toprule
pair & cosine obs/null & $p$ & top-3 overlap obs/null & $p$ & angle$^{\circ}$ obs/null & $p$ \\
\midrule
line--point   & 0.767 / 0.795 & 0.77 & 0.67 / 0.37 & 0.037 & 30 / 57 & 0.037 \\
line--saddle  & 0.813 / 0.730 & 1.00 & 0.50 / 0.35 & 0.049 & 45 / 59 & 0.049 \\
point--saddle & 0.792 / 0.756 & 0.92 & 0.67 / 0.36 & 0.020 & 30 / 58 & 0.020 \\
\bottomrule
\end{tabular}
\end{center}
The cosine (and correlation) of the full energy profiles did not differ from the random-quotient
null for any pair ($p\geq0.77$): all profiles place most of their energy on the single dominant
mode, which inflates cosine similarity equally in the observed and null cases, so cosine is not
discriminating here. The set- and subspace-based measures are: the three types shared more of their
top-3 modes ($0.50$--$0.67$ vs.\ $\approx0.36$) and their top-3 spans were closer ($30$--$45^{\circ}$
vs.\ $\approx58^{\circ}$) than under random quotients, significantly for every pair
($p=0.02$--$0.05$). Within a network the motif types thus occupy \emph{overlapping} dominant
mode-subspaces more than chance, rather than disjoint ones; the effect is visible in the top-mode
overlap and principal angles but not in the profile cosine.

\section{Cross-seed convergence of quotient modes}
\label{sec:crossseed}
\S\ref{sec:concentration}--\ref{sec:typeprofiles} are within-network. This section asks the
across-network question directly: do independently trained seeds develop the \emph{same} quotient
modes? Modes cannot be matched by index or by lifted direction (different unit spaces, no
correspondence), but the fifteen tasks are shared across seeds, so each mode is fingerprinted by a
seed-invariant \emph{task profile}: the $15$-vector $T[j,r]=\sum_{m\in r}(p_m\!\cdot\!M_j)^2$ of how
much each task's motifs load mode $j$ (its eigenvalue $\lvert\lambda\rvert$ is also seed-invariant).
For each seed we take the top-$m$ modes by aggregate motif energy and $L^1$-normalise their task
profiles. For every seed pair we Hungarian-match the two sets of modes by task-profile cosine and
record the mean matched cosine; the null applies a random permutation of the $15$-task axis to one
seed's profiles (destroying task identity, preserving profile shape) and rematches.

At $k=20$, top-$6$ modes: mean matched task-profile cosine $0.781$ versus task-shuffled null
$0.674$, with observed $>$ null in $44/45$ seed pairs (signed-rank $p<10^{-4}$). The effect is
robust across $k\in\{12,20,30\}$ and top-$m\in\{4,6,8\}$ (observed $0.75$--$0.80$ vs.\ null
$0.64$--$0.70$; $43$--$45$ of $45$ pairs; $p<10^{-4}$ throughout). The matched modes did \emph{not}
have similar eigenvalues (mean matched $\lvert\lambda\rvert$ gap $\approx7.9$): the convergence is
in task function, not spectral magnitude. Consistently, the tasks dominating the top modes across
all ten seeds were the delay / working-memory and context decision-making tasks
(\texttt{contextdelaydm1} dominant for a top mode in $10$ seeds, \texttt{delaydm1} in $9$,
\texttt{delaydm2} and \texttt{delayanti} in $7$, \texttt{contextdelaydm2} in $6$, \texttt{delaygo}
and \texttt{dmsgo} in $5$), and almost never the fast/reactive tasks (\texttt{fdgo}, \texttt{fdanti}
$\leq2$). Independently trained networks therefore converge on a shared repertoire of quotient
modes defined by which tasks recruit them --- a memory/decision-dominated set --- even though no
individual mode index or eigenvalue corresponds across seeds. This is the earned cross-network
version of the within-network concentration; it is a statement about task profiles, established by
matching, not about mode ordering.

\section{Eigenvalue (timescale) control}
\label{sec:lambda}
The motif directions are slow Jacobian modes ($\lvert\lambda\rvert$ near or above $1$). A reviewer
control: are containment, the quotient prediction $A_{\rm cond}$, and the top-quotient-mode energy
merely increasing functions of $\lvert\lambda\rvert$ (the implied timescale), so that ``the motif
lives in the connectivity subspace'' would just restate ``slow modes are captured''? For each
condition we took \emph{all} eigenvectors of the mean Jacobian $\bar J$ (not only the motif
direction) and recorded, per eigenvector, its containment $\lVert P_H v\rVert$, $A_{\rm cond}$
(against the same $r=3$ predicted-quotient basis), and its top-3 quotient-mode energy, together with
$\lvert\lambda\rvert$ and the timescale $\tau=-1/\log\lvert\lambda\rvert$. Pooling all $94{,}237$
eigenvectors over the ten networks ($k=20$):
\begin{center}\small
\begin{tabular}{lccc}
\toprule
 & vs $\lvert\lambda\rvert$: $r$ ($R^2$) & vs $\log\tau$: $r$ & shape across $\lvert\lambda\rvert$ \\
\midrule
containment      & $-0.05$ ($0.002$) & $-0.00$ & flat $0.27$ in bulk, rises to $0.68$ for $\lvert\lambda\rvert>1.05$ \\
$A_{\rm cond}$   & $+0.35$ ($0.12$)  & $+0.28$ & $0.35\to0.84\to0.94$ across the slow tail \\
top-3 energy     & $+0.07$ ($0.004$) & $+0.01$ & flat $0.07$ in bulk, rises to $0.73$ for $\lvert\lambda\rvert>1.05$ \\
\bottomrule
\end{tabular}
\end{center}
Globally none of the three is a strong linear function of $\lvert\lambda\rvert$ ($R^2\leq0.12$),
because slow modes ($\lvert\lambda\rvert\gtrsim0.95$) are only $\approx3\%$ of eigenvectors; but
within that slow tail all three rise steeply. The fair test is therefore a matched-$\lvert\lambda\rvert$
comparison: for each motif eigenvector, the quantity minus its mean over non-motif eigenvectors of
the \emph{same} $\lvert\lambda\rvert$ (window $\pm0.02$, per network), aggregated across seeds:
\begin{center}\small
\begin{tabular}{lccc}
\toprule
quantity & motif $-$ matched non-motif & seeds $>0$ & signed-rank $p$ \\
\midrule
containment      & $+0.022$ & $6/10$  & $0.43$ \\
$A_{\rm cond}$   & $+0.085$ & $9/10$  & $0.004$ \\
top-3 mode energy & $+0.199$ & $10/10$ & $0.002$ \\
\bottomrule
\end{tabular}
\end{center}
The results dissociate. \emph{Containment is a timescale phenomenon}: at matched
$\lvert\lambda\rvert$ the motif directions are contained no more than other Jacobian eigenvectors
($+0.02$, $p=0.43$) --- slow/near-marginal modes generally lie in the connectivity subspace, and the
motifs inherit this by being slow. (A naive residual against the near-flat global
$\lvert\lambda\rvert$--containment fit gives $+0.39$, $10/10$, $p=0.002$, but that overstates the
effect precisely because the global fit ignores the slow-tail rise; the matched control is the
correct test and it is null.) By contrast the quotient prediction and the concentration are
\emph{not} explained by timescale: at matched $\lvert\lambda\rvert$ the motif directions are still
better predicted by the reduced quotient ($A_{\rm cond}$ $+0.085$, $p=0.004$) and far more
concentrated on the top quotient modes ($+0.199$, $p=0.002$) than non-motif modes of the same
$\lvert\lambda\rvert$. Thus \S\ref{sec:containment} should be read as largely a statement about slow
modes, whereas \S\ref{sec:quotient} and \S\ref{sec:concentration} survive the timescale control.

\section{Activation robustness: released tanh networks}
\label{sec:tanh}
The main analyses use \texttt{softplus} networks. To test activation dependence we ran the same
pipeline on the three released \texttt{tanh} diagonal-init $256$-unit networks
(\texttt{LeakyRNN\_tanh\_diag\_256\_all\_0..2}): active units by the same $>10^{-3}$ variance
criterion (all $256$ units are active --- \texttt{tanh} does not sparsify as \texttt{softplus}
does), the signed-WL partition of the transpose-variant $W_{\rm rec}$ at the operating scale
($k=51$), fixed-point finding across the task epochs, and the mode concentration. We compare to the
three \emph{released} \texttt{softplus} seeds ($0$--$2$) so the sample size ($n=3$) and test are
matched.

\paragraph{Lumpability (Result 1): holds.} Eigenvector-leakage $z$ of $W_{\rm rec}$ onto the
group-constant subspace vs.\ size-matched random partitions was strongly negative in all three
\texttt{tanh} networks: transpose-orientation leakage $z=-17.7,\,-17.1,\,-19.8$ (invariance
$z=-27,\,-23,\,-28$; original-orientation leakage $z=-8.0,\,-6.3,\,-8.4$). These are in the same
range as the \texttt{softplus} networks (\S\ref{sec:lump}). The signed connectivity partition is
lumpable under \texttt{tanh} as well.

\paragraph{Motif inventory differs.} The \texttt{tanh} fixed-point inventory is
\emph{point-attractor dominated} ($18/11/6$, $18/5/0$, $22/5/2$ point/line/saddle per seed; $58$
points, $21$ lines, $8$ saddles over the three nets), unlike the line-attractor--dominated
\texttt{softplus} inventory (\S\ref{sec:motifs}).

\paragraph{Concentration (Result 4): does not hold.} Motif energy did \emph{not} concentrate on the
connectivity quotient's modes beyond a random quotient in the \texttt{tanh} networks, in contrast
to \texttt{softplus} on the matched three released seeds:
\begin{center}\small
\begin{tabular}{llccc}
\toprule
activation & $k$ & PR conn / rand ($z$) & top-3 conn / rand & median per-net perm.\ $p$ \\
\midrule
\texttt{softplus} ($0$--$2$) & 12 & 2.85 / 6.90 ($-2.8$) & 0.785 / 0.551 & 0.003 \\
                             & 20 & 3.67 / 10.64 ($-3.2$) & 0.797 / 0.426 & 0.003 \\
                             & 30 & 5.04 / 15.36 ($-3.6$) & 0.670 / 0.337 & 0.003 \\
\texttt{tanh} ($0$--$2$)     & 12 & 5.88 / 6.29 ($-0.3$) & 0.587 / 0.586 & 0.37 \\
                             & 20 & 10.36 / 10.55 ($+0.7$) & 0.437 / 0.428 & 0.71 \\
                             & 30 & 19.23 / 14.22 ($+1.4$) & 0.264 / 0.350 & 0.92 \\
\bottomrule
\end{tabular}
\end{center}
At matched $n$ and with the identical test, the \texttt{softplus} motifs concentrated far below
their own random-quotient null (per-network $z\approx-3$; top-3 energy $\approx0.8$ vs.\ $0.4$),
while the \texttt{tanh} motifs did not (per-network $z\approx0$; top-3 energy at the random level).
By type at $k=20$ the only hint of \texttt{tanh} concentration was in the few saddles
($z\approx-1.1$, $n=6$); the dominant point attractors did not concentrate ($z\approx+0.6$). Thus
the lumpability of the connectivity partition is activation-general, but the localisation of
fixed-point motifs on the quotient's dominant modes is specific to the (line-attractor--dominated)
\texttt{softplus} networks and does not transfer to the (point-attractor--dominated) \texttt{tanh}
networks. The three tanh networks are a small sample and this is reported as a robustness check,
not a seed-level significance test.

\section{Cross-task participation reuse and its control (negative)}
\label{sec:reuse}
For each motif a participation vector over the groups (summed squared loading) was formed; the
reuse statistic $R$ is the mean same-computation correlation minus the mean different-computation
correlation. Against a size-matched random partition (each minus its own baseline), across the
ten networks:
\begin{center}\small
\begin{tabular}{rccccc}
\toprule
$k$ & $R$ (conn) & $R$ (random) & $R$ (activity) & conn $>$ rand & $p$ \\
\midrule
8  & 0.127 & 0.195 & 0.191 & 0/10 & 0.002 \\
12 & 0.123 & 0.193 & 0.180 & 2/10 & 0.010 \\
20 & 0.147 & 0.197 & 0.180 & 3/10 & 0.193 \\
30 & 0.158 & 0.194 & 0.195 & 1/10 & 0.020 \\
50 & 0.194 & 0.197 & 0.201 & 6/10 & 0.695 \\
\bottomrule
\end{tabular}
\end{center}
The connectivity participation reuse did not exceed the size-matched random partition at any $k$
(median $R_{\rm conn}<R_{\rm random}$ throughout; connectivity above random in $0$--$6$ of $10$
networks depending on $k$, never with $p<0.05$ in its favour).

\subsection{Relation to the concentration result (\S\ref{sec:concentration})}
Because connectivity motif energy concentrates on few quotient modes (\S\ref{sec:concentration}),
one might expect the reuse deficit here to be a mechanical consequence: concentrated
participation vectors are supported on fewer groups, which could compress the variance of
between-motif correlations and lower $R$. We tested this directly at $k=20$ using the reuse
statistic above (cross-task pairs, group-projected profiles). The connectivity participation
vectors were indeed more concentrated than random --- median effective group support $4.9$--$7.7$
versus $6.5$--$10.5$, lower in $10/10$ networks --- but the reuse deficit was null-sized
($\overline{R_{\rm conn}-R_{\rm random}}=-0.013$) and did \emph{not} track concentration across
seeds (Spearman $\rho=-0.03$, $p=0.93$ between the per-seed support gap and reuse gap). The two
observations are therefore reported as separate findings: concentration is a robust positive
against random quotients (\S\ref{sec:concentration}); the reuse deficit is a small, essentially
null-sized negative (this section) that concentration does not explain.

\section{Trajectory decoding from the quotient coordinates}
\label{sec:decode}
Does the $k$-dimensional connectivity quotient carry the computation, or does group-averaging
discard the task-relevant state? For each network we ran the fifteen tasks (noise off) and, at
every retained epoch, took the instantaneous hidden state at four sampled times. We decoded two
trajectory variables --- task identity ($15$-way) and computational phase (fixation / stimulus /
delay / response, pooled across tasks) --- with balanced multinomial logistic regression
($5$-fold CV) from four matched feature sets: the full active state ($\approx210$ dims, the
ceiling), the WL role-group means ($k$ dims), a size-matched random-partition's group means
($k$), and a random Gaussian projection to $k$ dims. Mean CV accuracy over the ten networks
($k=20$):
\begin{center}\small
\begin{tabular}{lcccccl}
\toprule
variable & full & WL quotient & random partition & random projection & chance & WL vs rand-proj \\
\midrule
task identity     & 1.000 & 0.727 & 0.750 & 0.733 & 0.223 & $5/10$, $p=0.62$ \\
computational phase & 0.997 & 0.889 & 0.812 & 0.791 & 0.311 & $10/10$, $p=0.002$ \\
\bottomrule
\end{tabular}
\end{center}
The two variables dissociate. Computational phase was decoded from the $k$ WL quotient
coordinates better than from either $k$-dimensional random baseline ($10/10$ networks,
signed-rank $p=0.002$ against both), i.e.\ the coarse connectivity coordinates preferentially
retain the dynamical stage of the trajectory. Task identity was not: the WL quotient decoded it no
better than a random projection of the same dimension ($5/10$, $p=0.62$). This is consistent with
the quotient capturing the slow dynamical state (which carries the phase) rather than the
input-driven task label (which the full state decodes at ceiling but the slow subspace does not
privilege).

\section{Causal block-model resample (negative)}
\label{sec:causal}
The results above are correlational: they measure how far the quotient coordinates track the
trained dynamics. This is a causal test of whether preserving the quotient $Q_{\rm struct}$ is
sufficient for behaviour. We resampled the trained recurrent matrix under two surrogates sharing
the same global weight multiset but differing in whether the block structure is kept, injected each
back into the network with all other weights fixed and no retraining, and measured mean performance
on the fifteen tasks:
\begin{itemize}\itemsep0pt
\item \emph{$Q$-preserving}: permute residuals within each $(g_i,g_j)$ role block. This preserves
every block mean, hence $Q_{\rm struct}=D^{-1}H^{T}W_{a}H$ exactly, while destroying within-block
wiring.
\item \emph{partition-agnostic}: permute all active--active entries globally. Same weight multiset,
but block means (and $Q_{\rm struct}$) destroyed.
\end{itemize}
Means over the ten $256$-unit networks ($k=20$, five surrogate draws each); performance is the mean
fraction correct over the fifteen tasks, the spectral error is the mean $\lvert\Delta\lambda\rvert$
of the top-ten $\lvert\lambda\rvert$ of the active recurrent matrix relative to the trained one.
\begin{center}\small
\begin{tabular}{lccc}
\toprule
 & trained & $Q$-preserving & partition-agnostic \\
\midrule
task performance                       & 0.944 & 0.044 & 0.042 \\
top-10 $\lvert\lambda\rvert$ error     & ---   & 6.16  & 9.28 \\
\bottomrule
\end{tabular}
\end{center}
Both surrogates collapsed performance from the trained ceiling ($0.944$) to near the fixation
floor ($\approx0.04$), and the $Q$-preserving surrogate was \emph{not} better than the
partition-agnostic one ($Q$-preserving above partition-agnostic in $6/10$ networks, signed-rank
$p=0.62$). Preserving the coarse $20$-group quotient (a few hundred block means out of
$\sim\!4\times10^{4}$ weights) is thus not causally sufficient for the computation: the fine
within-block wiring carries the behaviour. However, the $Q$-preserving surrogate did preserve the
trained matrix's leading eigenvalues markedly better than the partition-agnostic shuffle (spectral
error $6.16$ vs.\ $9.28$; $Q$-preserving lower in $10/10$ networks, $p=0.002$). The two readouts
dissociate: preserving $Q_{\rm struct}$ retains the slow eigenspectrum --- the structure the
correlational results (\S\ref{sec:lump}--\ref{sec:decode}) concern --- but not enough of the fine
wiring to retain behaviour. This bounds the quotient claim: the connectivity coordinates are a real
property of the slow dynamical subspace, not a sufficient causal substitute for the trained weights.

\section{Statistical conventions}
\label{sec:statconv}
Cross-network claims treated the seed as the unit of replication; effects were tested across
seeds with a two-sided Wilcoxon signed-rank test ($p=0.002$ is the floor for $n=10$).
Within-network nulls preserved group sizes and permuted unit labels (or, for similarity
analyses, unit identities); each surrogate or size-matched partition was scored against its own
recomputed null. Standardised scores are relative to the null mean and standard deviation,
oriented positive-for-more-structure. Permutation nulls used $1000$ draws (a first pass at $200$
gave the same signs and comparable magnitudes).

\paragraph{Multiple comparisons.} The report contains one confirmatory family and a set of
control/exploratory tests. The confirmatory family is the four backbone results, each as a single
across-seed signed-rank test at the operating scale: (R1) lumpability leakage, (R2) quotient
prediction $>$ random, (R3) rotation ring $>$ matched, (R4) top-3 concentration $>$ random. Four
tests: at Bonferroni $\alpha=0.05/4=0.0125$ all four remain significant (R1, R2, R4 at the
$p=0.002$ floor; R3 by the $104/109$ binomial). The $k$-sweep entries for a given result are the
same hypothesis re-tested at neighbouring scales, not independent members of the family, and are
reported as a robustness curve rather than counted separately. The negative controls
(eigenvalue-matched, task-accessibility, cross-task reuse), the strong-null comparisons
($W_{\rm rec}$-eigenbasis, input-weight span), the surrogate mechanism, and the accuracy
correlations are diagnostic and are reported at face value without family correction; none is
used to support a positive claim.

\section{Data and code provenance}
Analyses were run under the \texttt{flexsym} package. The result tables and JSON summaries are:
\begin{itemize}\itemsep0pt
\item \texttt{lumpability\_ksweep.csv}, \texttt{lumpability.csv} (\S\ref{sec:lump});
\item \texttt{training\_control\_all\_sizes.csv} (\S\ref{sec:surrogate});
\item \texttt{quotient\_prediction.json}, \texttt{quotient\_controls.json} (\S\ref{sec:quotient});
\item \texttt{quotient\_eigmatched\_control.json} (\S\ref{sec:eigmatched});
\item \texttt{subspace\_role\_vs\_activity.json} (\S\ref{sec:containment});
\item \texttt{rotation\_readout\_control.csv} (\S\ref{sec:rotation});
\item \texttt{quotient\_mode\_concentration.json} (\S\ref{sec:concentration});
\item \texttt{task\_accessibility.csv} (\S\ref{sec:access});
\item \texttt{reuse\_random\_partition\_control.json} (\S\ref{sec:reuse});
\item \texttt{motif\_taxonomy\_all\_nets.csv} (\S\ref{sec:motifs}).
\end{itemize}
The uncorrected first-batch models and superseded outputs are documented in
\texttt{legacy\_uncorrected\_seeds\_report}.

\appendix
\section{Per-seed tables}
All tables are the ten $256$-unit networks (seeds $0$--$2$ released, $3$--$9$ retrained on the
corrected objective) unless noted. ``score'' denotes a signed standardised deviation from a
size-matched null (positive-for-more-structure).

\subsection{Lumpability, per seed (leakage and invariance scores at $k=8,20,50$; \S\ref{sec:lump})}
\begin{center}\small
\begin{tabular}{llrrrrrr}
\toprule
& & \multicolumn{3}{c}{leak score ($W^{T}$)} & \multicolumn{3}{c}{invariance score}\\
\cmidrule(lr){3-5}\cmidrule(lr){6-8}
seed & model & $k{=}8$ & $k{=}20$ & $k{=}50$ & $k{=}8$ & $k{=}20$ & $k{=}50$ \\
\midrule
0 & released & +32.1 & +27.4 & +16.6 & $-2.7$ & +1.8 & +17.6 \\
1 & released & +51.2 & +27.3 & +19.2 & $-12.1$ & $-1.6$ & +11.1 \\
2 & released & +28.8 & +27.3 & +15.8 & +1.1 & +3.8 & +21.4 \\
3 & retrained & +14.0 & +17.9 & +17.9 & +1.1 & +5.0 & +19.8 \\
4 & retrained & +25.1 & +28.9 & +19.8 & +1.1 & +9.4 & +17.8 \\
5 & retrained & +27.5 & +28.1 & +20.9 & +1.1 & +9.6 & +15.3 \\
6 & retrained & +37.7 & +29.9 & +16.8 & $-0.4$ & +7.8 & +18.8 \\
7 & retrained & +28.0 & +31.2 & +19.5 & +0.9 & +10.4 & +15.8 \\
8 & retrained & +21.0 & +23.0 & +20.3 & +1.1 & +8.1 & +21.9 \\
9 & retrained & +20.3 & +30.5 & +20.8 & +1.2 & +6.2 & +16.5 \\
\bottomrule
\end{tabular}
\end{center}

\subsection{Degree/training surrogate control, per network (\S\ref{sec:surrogate})}
$256$-unit seeds $0$--$9$ (corrected); $128/512/1024$-unit seeds $0$--$2$ (released).
\begin{center}\footnotesize\setlength{\tabcolsep}{4pt}
\resizebox{\textwidth}{!}{
\begin{tabular}{llrrrrrrrr}
\toprule
& & \multicolumn{2}{c}{trained} & \multicolumn{2}{c}{row shuffle} & \multicolumn{2}{c}{full shuffle} & \multicolumn{2}{c}{Gaussian}\\
\cmidrule(lr){3-4}\cmidrule(lr){5-6}\cmidrule(lr){7-8}\cmidrule(lr){9-10}
size & seed & eq & $R^{2}$ & eq & $R^{2}$ & eq & $R^{2}$ & eq & $R^{2}$ \\
\midrule
128 & 0 & +3.1 & 0.753 & +12.1 & 0.778 & +0.7 & 0.385 & +3.3 & 0.338 \\
128 & 1 & +5.5 & 0.816 & +11.5 & 0.799 & +2.8 & 0.428 & +1.7 & 0.384 \\
128 & 2 & +4.4 & 0.825 & +12.6 & 0.786 & +3.1 & 0.426 & +3.7 & 0.278 \\
256 & 0 & +4.2 & 0.831 & +14.8 & 0.827 & +1.5 & 0.455 & +2.4 & 0.333 \\
256 & 1 & +5.7 & 0.890 & +10.5 & 0.577 & +2.6 & 0.355 & +3.5 & 0.245 \\
256 & 2 & +7.5 & 0.834 & +16.0 & 0.753 & +0.5 & 0.411 & +2.1 & 0.279 \\
256 & 3 & +8.7 & 0.727 & +13.7 & 0.681 & +4.5 & 0.433 & +2.8 & 0.260 \\
256 & 4 & +7.2 & 0.847 & +8.1 & 0.677 & +3.1 & 0.352 & +2.5 & 0.301 \\
256 & 5 & +9.0 & 0.673 & +15.1 & 0.643 & +5.5 & 0.476 & +1.8 & 0.254 \\
256 & 6 & +7.4 & 0.756 & +10.3 & 0.679 & +4.3 & 0.368 & +2.4 & 0.227 \\
256 & 7 & +7.0 & 0.750 & +11.1 & 0.722 & +4.4 & 0.434 & +1.6 & 0.279 \\
256 & 8 & +6.9 & 0.746 & +10.3 & 0.725 & +2.9 & 0.406 & $-0.5$ & 0.305 \\
256 & 9 & +10.0 & 0.787 & +12.7 & 0.649 & +2.1 & 0.459 & +2.6 & 0.268 \\
512 & 0 & +6.0 & 0.839 & +17.3 & 0.748 & +1.9 & 0.288 & +4.3 & 0.274 \\
512 & 1 & +1.7 & 0.844 & +18.0 & 0.777 & +2.4 & 0.336 & +2.8 & 0.294 \\
512 & 2 & +5.3 & 0.857 & +19.5 & 0.818 & +3.1 & 0.398 & +1.7 & 0.296 \\
1024 & 0 & $-4.2$ & 0.835 & +17.0 & 0.769 & +1.8 & 0.348 & +1.9 & 0.218 \\
1024 & 1 & $-6.2$ & 0.851 & +18.8 & 0.778 & +1.2 & 0.336 & +3.6 & 0.219 \\
1024 & 2 & +0.6 & 0.878 & +20.3 & 0.843 & +3.0 & 0.398 & +2.0 & 0.275 \\
\bottomrule
\end{tabular}}
\end{center}

\subsection{Quotient prediction, per seed ($k=20$, $r=3$; \S\ref{sec:quotient})}
\begin{center}\small
\begin{tabular}{llrrrrrr}
\toprule
seed & model & $A_{\rm cond}$ & rand & activity & $z$ & $|\Delta\lambda|$ conn & rand \\
\midrule
0 & released & 0.889 & 0.642 & 0.703 & +4.02 & 0.047 & 0.193 \\
1 & released & 0.896 & 0.627 & 0.626 & +4.19 & 0.031 & 0.195 \\
2 & released & 0.846 & 0.637 & 0.646 & +3.65 & 0.043 & 0.180 \\
3 & retrained & 0.807 & 0.612 & 0.603 & +3.10 & 0.056 & 0.187 \\
4 & retrained & 0.821 & 0.584 & 0.658 & +3.81 & 0.038 & 0.188 \\
5 & retrained & 0.872 & 0.603 & 0.625 & +4.27 & 0.046 & 0.193 \\
6 & retrained & 0.774 & 0.614 & 0.547 & +2.41 & 0.042 & 0.189 \\
7 & retrained & 0.791 & 0.589 & 0.621 & +3.25 & 0.048 & 0.198 \\
8 & retrained & 0.811 & 0.610 & 0.584 & +3.03 & 0.046 & 0.188 \\
9 & retrained & 0.758 & 0.613 & 0.685 & +2.26 & 0.076 & 0.202 \\
\bottomrule
\end{tabular}
\end{center}

\subsection{Motif-direction containment, per seed ($k=20$; \S\ref{sec:containment})}
\begin{center}\small
\begin{tabular}{llrrrr}
\toprule
seed & model & \#motifs & $\Delta_{\rm role}$ & $\Delta_{\rm activity}$ & $z$ (role$-$act) \\
\midrule
0 & released & 47 & 0.329 & $-0.000$ & +8.84 \\
1 & released & 33 & 0.322 & 0.041 & +8.38 \\
2 & released & 48 & 0.285 & 0.067 & +6.33 \\
3 & retrained & 44 & 0.210 & 0.022 & +5.09 \\
4 & retrained & 43 & 0.243 & $-0.004$ & +7.30 \\
5 & retrained & 42 & 0.205 & $-0.007$ & +5.42 \\
6 & retrained & 48 & 0.335 & 0.016 & +7.31 \\
7 & retrained & 43 & 0.211 & 0.047 & +5.18 \\
8 & retrained & 45 & 0.280 & $-0.010$ & +8.45 \\
9 & retrained & 47 & 0.279 & 0.007 & +7.57 \\
\bottomrule
\end{tabular}
\end{center}

\subsection{Rotation into the readout, per seed (\S\ref{sec:rotation})}
Median readout potency over each network's delayed-response tasks.
\begin{center}\small
\begin{tabular}{llrrrrr}
\toprule
seed & model & tasks & ring (delay) & ring (resp.) & matched & random \\
\midrule
0 & released & 11 & 0.0000 & 0.1014 & 0.0058 & 0.0098 \\
1 & released & 11 & 0.0001 & 0.1007 & 0.0018 & 0.0081 \\
2 & released & 11 & 0.0000 & 0.0631 & 0.0041 & 0.0096 \\
3 & retrained & 11 & 0.0009 & 0.0288 & 0.0006 & 0.0079 \\
4 & retrained & 10 & 0.0007 & 0.0196 & 0.0020 & 0.0086 \\
5 & retrained & 11 & 0.0003 & 0.0293 & 0.0004 & 0.0088 \\
6 & retrained & 11 & 0.0002 & 0.0291 & 0.0045 & 0.0080 \\
7 & retrained & 11 & 0.0127 & 0.0347 & 0.0004 & 0.0079 \\
8 & retrained & 11 & 0.0006 & 0.0427 & 0.0009 & 0.0082 \\
9 & retrained & 11 & 0.0008 & 0.0197 & 0.0015 & 0.0081 \\
\bottomrule
\end{tabular}
\end{center}
Ring-response $>$ matched (per task): all tasks for seeds $0,1,2,3,5,6,8,9$; $9/10$ (seed $4$);
$7/11$ (seed $7$).

\subsection{Mode concentration, per seed (top-3 energy and PR; \S\ref{sec:concentration})}
\begin{center}\footnotesize\setlength{\tabcolsep}{4pt}
\resizebox{\textwidth}{!}{
\begin{tabular}{llrrrrrrrr}
\toprule
& & \multicolumn{2}{c}{PR $k{=}12$} & \multicolumn{2}{c}{PR $k{=}20$} & \multicolumn{2}{c}{PR $k{=}30$} & \multicolumn{2}{c}{top-3 ($k{=}20$)}\\
\cmidrule(lr){3-4}\cmidrule(lr){5-6}\cmidrule(lr){7-8}\cmidrule(lr){9-10}
seed & model & conn & rand & conn & rand & conn & rand & conn & rand \\
\midrule
0 & released & 4.04 & 6.60 & 4.13 & 10.27 & 5.37 & 13.88 & 0.80 & 0.44 \\
1 & released & 2.85 & 7.83 & 3.37 & 11.45 & 4.45 & 15.89 & 0.80 & 0.41 \\
2 & released & 2.52 & 6.90 & 3.67 & 10.64 & 5.04 & 15.36 & 0.67 & 0.43 \\
3 & retrained & 2.17 & 7.27 & 4.60 & 10.81 & 7.81 & 15.26 & 0.60 & 0.42 \\
4 & retrained & 4.23 & 6.45 & 6.59 & 10.46 & 9.26 & 13.77 & 0.60 & 0.43 \\
5 & retrained & 2.40 & 6.79 & 5.15 & 9.22 & 7.01 & 13.16 & 0.63 & 0.47 \\
6 & retrained & 2.18 & 5.40 & 6.30 & 8.66 & 6.42 & 13.01 & 0.54 & 0.49 \\
7 & retrained & 4.03 & 7.60 & 6.13 & 9.35 & 8.02 & 14.14 & 0.64 & 0.46 \\
8 & retrained & 2.77 & 7.16 & 4.40 & 12.08 & 6.65 & 16.88 & 0.70 & 0.39 \\
9 & retrained & 2.66 & 7.59 & 5.45 & 11.53 & 6.54 & 16.18 & 0.64 & 0.40 \\
\bottomrule
\end{tabular}}
\end{center}
Per-seed task accuracy (for the accuracy correlation, \S\ref{sec:concentration}): seeds $0$--$9$
$=0.991,0.974,0.987,0.975,0.976,0.974,0.974,0.973,0.974,0.975$.

\subsection{Eigenvalue-matched non-motif control, per seed ($k=20$; \S\ref{sec:eigmatched})}
\begin{center}\small
\begin{tabular}{llrrrr}
\toprule
seed & model & \#cond & motif & matched & win \\
\midrule
0 & released & 47 & 0.889 & 0.914 & 15/47 \\
1 & released & 33 & 0.896 & 0.744 & 18/33 \\
2 & released & 48 & 0.846 & 0.806 & 24/48 \\
3 & retrained & 44 & 0.807 & 0.709 & 26/44 \\
4 & retrained & 43 & 0.821 & 0.765 & 23/43 \\
5 & retrained & 42 & 0.872 & 0.861 & 13/42 \\
6 & retrained & 48 & 0.774 & 0.792 & 23/48 \\
7 & retrained & 43 & 0.791 & 0.820 & 17/43 \\
8 & retrained & 45 & 0.811 & 0.746 & 25/45 \\
9 & retrained & 47 & 0.758 & 0.752 & 20/47 \\
\bottomrule
\end{tabular}
\end{center}

\subsection{Cross-task participation reuse, per seed ($k=20$; \S\ref{sec:reuse})}
\begin{center}\small
\begin{tabular}{llrrrr}
\toprule
seed & model & $R$ conn & $R$ random & $R$ activity & $\Delta$ \\
\midrule
0 & released & 0.305 & 0.256 & 0.272 & +0.048 \\
1 & released & 0.136 & 0.141 & 0.155 & $-0.005$ \\
2 & released & 0.129 & 0.176 & 0.193 & $-0.046$ \\
3 & retrained & 0.128 & 0.185 & 0.115 & $-0.058$ \\
4 & retrained & 0.141 & 0.221 & 0.148 & $-0.080$ \\
5 & retrained & 0.140 & 0.205 & 0.167 & $-0.065$ \\
6 & retrained & 0.202 & 0.191 & 0.167 & +0.011 \\
7 & retrained & 0.238 & 0.202 & 0.193 & +0.036 \\
8 & retrained & 0.153 & 0.189 & 0.217 & $-0.036$ \\
9 & retrained & 0.200 & 0.206 & 0.225 & $-0.006$ \\
\bottomrule
\end{tabular}
\end{center}
$\Delta>0$ in $2/10$ networks (seeds $6,7$; and seed $0$).

\subsection{Task accessibility, per seed (seeds $0$--$2$; \S\ref{sec:access})}
Number of conditions with motif mode $>$ matched mode, per measure.
\begin{center}\footnotesize\setlength{\tabcolsep}{4pt}
\resizebox{\textwidth}{!}{
\begin{tabular}{lrrrrrrrr}
\toprule
seed & \#cond & input & displ.\ & traj.\ var & readout & persist.\ & controll.\ & observ.\ \\
\midrule
0 & 47 & 13/47 & 19/47 & 33/47 & 22/47 & 32/47 & 19/47 & 23/47 \\
1 & 33 & 9/33 & 13/33 & 19/33 & 18/33 & 12/33 & 5/33 & 18/33 \\
2 & 48 & 23/48 & 30/48 & 34/48 & 13/48 & 31/48 & 21/48 & 14/48 \\
\bottomrule
\end{tabular}}
\end{center}

\subsection{Fixed-point taxonomy, per seed (corrected models; \S\ref{sec:motifs})}
\begin{center}\footnotesize\setlength{\tabcolsep}{4pt}
\resizebox{\textwidth}{!}{
\begin{tabular}{lrrrrrrrr}
\toprule
& \multicolumn{3}{c}{gate (matched / total)} & \multicolumn{5}{c}{inventory (\#epochs)}\\
\cmidrule(lr){2-4}\cmidrule(lr){5-9}
seed & mem$\to$line & dec$\to$saddle & sens$\to$point & line & point & saddle & mixed & rot. \\
\midrule
0 & 16/16 & 14/14 & 4/4 & 26 & 12 & 9 & 2 & 0 \\
1 & 8/16 & 3/14 & 1/4 & 16 & 14 & 3 & 16 & 0 \\
2 & 16/16 & 13/14 & 4/4 & 25 & 15 & 8 & 1 & 0 \\
3 & 10/16 & 12/14 & 1/4 & 27 & 10 & 7 & 5 & 0 \\
4 & 13/16 & 13/14 & 1/4 & 27 & 6 & 10 & 6 & 0 \\
5 & 13/16 & 10/14 & 0/4 & 26 & 8 & 8 & 7 & 0 \\
6 & 16/16 & 7/14 & 1/4 & 32 & 11 & 5 & 1 & 0 \\
7 & 12/16 & 12/14 & 1/4 & 27 & 9 & 7 & 6 & 0 \\
8 & 12/16 & 13/14 & 0/4 & 27 & 9 & 9 & 4 & 0 \\
9 & 14/16 & 12/14 & 0/4 & 30 & 7 & 10 & 2 & 0 \\
\bottomrule
\end{tabular}}
\end{center}

\end{document}
